\documentclass[../../main.tex]{subfiles}
\begin{document}

\begin{flushright}
\footnotesize\textit{【自動文字起こし・要確認】}
\end{flushright}

\noindent
{\bf [解]}
\paragraph{(1)}

$[3^m]$は$3^m$で割り切れるが，$3^{m+1}$では割り切れない（$\cdots\diamondsuit$）ことを帰納的に示す．$m=0$の時，$3^0=1$，$[1]=1$だから$\diamondsuit$は成立するので，以下$m=k\in\mathbb Z_{\ge0}$での成立を仮定し，$m=k+1$でも$\diamondsuit$が成立することを示す．$A=10^{3^k}$とおく．
\begin{align}
[3^{k+1}]=\dfrac{10^{3^{k+1}}-1}{9}=\dfrac{A^3-1}{9}=\dfrac{A-1}{9}(A^2+A+1)\quad\cdots①
\end{align}
ここで，$\dfrac{A-1}{9}=[3^k]$だから仮定より，$\dfrac{A-1}{9}$は$3^k$で割り切れ$3^{k+1}$で割り切れない．一方，
\begin{align}
A^2+A+1&\equiv1+1+1\equiv0\pmod3\\
A^2+A+1&\equiv3\not\equiv0\pmod9
\end{align}
（$\because10\equiv1\pmod9$より$A\equiv1\pmod9$）だから，$A^2+A+1$は$3$で割り切れるが$9$で割り切れない．以上及び①から，$[3^{k+1}]$は$3^{k+1}$で割り切れ，$3^{k+2}$で割り切れないので，$m=k+1$でも$\diamondsuit$は成立．よって題意は示された．$\blacksquare$

\paragraph{(2)}

$P$: 「$n$が$27$で割り切れる」，$Q$: 「$[n]$が$27$で割り切れる」とおく．

$1^\circ$ 十分性

$P$が成立する時，$3$と互いに素な自然数$t$，及び$k\in\mathbb N_{\ge3}$を用いて，$n=t\cdot3^k$とかける．$A=10^{3^k}$として，
\begin{align}
[n]=\dfrac{A^t-1}{9}=\dfrac{A-1}{9}(A^{t-1}+\cdots+A+1)
\end{align}
で，(1)から$\dfrac{A-1}{9}=[3^k]$は$3^k$で割り切れる．$k\ge3$からつまり$[3^k]$は$27$で割り切れるので，$[n]$も$27$で割り切れる．たしかに$P\to Q$は成立．

$2^\circ$ 必要性

$n$は$3$と互いに素な自然数$t$，及び$k\in\mathbb Z_{\ge0}$を用いて，$n=t\cdot3^k$とかける．$Q$が成立する時，
\begin{align}
[n]=\dfrac{A^t-1}{9}=\dfrac{A-1}{9}(A^{t-1}+\cdots+A+1)
\end{align}
において，
\begin{align}
A^{t-1}+\cdots+A+1\equiv\underbrace{1+1+\cdots+1}_{t個}\equiv t\not\equiv0\pmod3\quad(\because t\perp3)
\end{align}
だから，$\dfrac{A-1}{9}$が$27$で割り切れる．(1)からしたがって，$k\ge3$となり，この時$n$が$27$で割り切れるから，たしかに$Q\to P$も成立する．

以上から，$P\iff Q$となり，題意は示された．$\blacksquare$

\bigskip

\end{document}
